\customSection[sec:analysis]{Аналіз предметної області}{0em}

\customSubSection[sec:analysis_two_level_section]{Оформлення коду}{0em}

Оформлення коду є важливим аспектом при написанні технічної документації. Нижче наведено приклади оформлення коду для різних мов програмування.

% ----------------------
% Example 1: TypeScript
% ----------------------
TypeScript:
\begin{lstlisting}[style=nocodeframe, language=TypeScript]
interface User {
    name: string;
    age: number;
}

const user: User = {
    name: "Ivan",
    age: 17
};

function greet(u: User): void {
    console.log(`Hello, ${u.name}`);
}

greet(user);
\end{lstlisting}

% ----------------------
% Example 2: Arduino C++
% ----------------------
ArduinoCpp:
\begin{lstlisting}[style=nocodeframe, language=ArduinoCpp]
#include <Arduino.h>

void setup() {
    Serial.begin(9600);
}

void loop() {
    Serial.println("Hello Arduino");
}
\end{lstlisting}

Lorem Ipsum — це стандартний заповнювач тексту в поліграфії та верстці, який використовується з 1500-х років. Він допомагає демонструвати макети без використання реального змісту.

\customSubSection[sec:analysis_two_level_section]{Оформлення рисунків}{1em}

Оформлення рисунків є важливим аспектом при написанні технічної документації. Нижче наведено приклад оформлення рисунка.

\setfigure[fig:intro_example2]{0.8\textwidth}{assets/img}{Тестовий рисунок}

На \textbf{рис.~\ref{fig:intro_example2}} показано приклад зображення, яке ілюструє відповідний приклад.

\customSubSection[sec:analysis_third_level_section]{Оформлення таблиць}{1em}

\begin{spacing}{1.3}
    \noindent
    \begin{xltabular}{\textwidth}{| c | >{\raggedright\arraybackslash}X | c |}
        \caption{Список дисциплін} \label{tab:disciplines} \\
        \hline
        \textbf{№} & \textbf{Назва дисципліни} & \textbf{К-ть кредитів ЄКТС} \\ \hline
        \endfirsthead

        \multicolumn{3}{r}{\textit{}} \\ \hline
        \textbf{№} & \textbf{Назва дисципліни} & \textbf{К-ть кредитів ЄКТС} \\ \hline
        \endhead

        1 & Методи та засоби інженерії програмного забезпечення & 220 \\ \hline
        1 & Методи та засоби інженерії програмного забезпечення & 220 \\ \hline
        1 & Методи та засоби інженерії програмного забезпечення & 220 \\ \hline
        1 & Методи та засоби інженерії програмного забезпечення & 220 \\ \hline
        1 & Методи та засоби інженерії програмного забезпечення & 220 \\ \hline
    \end{xltabular}
\end{spacing}

\customSubSection[sec:analysis_third_level_section]{Міждокументні посилання}{1em}

Як зазначено в підрозділі \ref{sec:analysis_two_level_section}, аналіз предметної області є критично важливим.

На \textbf{рис.~\ref{fig:intro_example2}} продемонстровано приклад зображення, яке ілюструє відповідний приклад.

Тема Arduino розглядається детально у багатьох джерелах \citecustom{1}, \citecustom{3}, а практичні проекти описані у \citecustom{4}.

Дані дисциплін наведені у \tabref{tab:disciplines}.

\begin{spacing}{1.3}
    \noindent
    \begin{xltabular}{\textwidth}{| c | >{\raggedright\arraybackslash}X | c |}
        \caption{Ще одна тестова таблиця} \label{tab:test_table} \\
        \hline
        \textbf{№} & \textbf{Тест} & \textbf{Тест} \\ \hline
        \endfirsthead

        \multicolumn{3}{r}{\textit{}} \\ \hline
        \textbf{№} & \textbf{Тест} & \textbf{Тест} \\ \hline
        \endhead

        1 & Методи та засоби інженерії програмного забезпечення & 220 \\ \hline
    \end{xltabular}
\end{spacing}